% ============================================================================
% Section V: Software Architecture
% ============================================================================

\section{Software Architecture}
\label{sec:architecture}

\subsection{Design Principles}

\gtopt{} is designed around four core principles:

\begin{enumerate}
  \item \textbf{Performance}: C++26 with flat-map-based sparse matrices
    for fast LP assembly.  The \texttt{boost::container::flat\_map}
    (aliased as \texttt{gtopt::flat\_map}) provides 10--27$\times$ faster
    iteration than \texttt{std::map} in LP coefficient assembly.
  \item \textbf{Modularity}: Each power system component (generator,
    demand, line, battery, reservoir, etc.) has its own LP assembly class
    (\texttt{*\_lp.cpp}) that independently adds variables, constraints,
    and objective terms to the shared \texttt{LinearProblem}.
  \item \textbf{Extensibility}: New components and solver backends can be
    added without modifying the core LP engine.  Solver backends are
    loaded as dynamic plugins (\texttt{libgtopt\_solver\_*.so}).
  \item \textbf{Interoperability}: JSON configuration with native Apache
    Arrow/Parquet I/O enables seamless integration with Python, R, and
    other data science ecosystems.
\end{enumerate}

\subsection{System Architecture}

Fig.~\ref{fig:arch_layers} illustrates the layered architecture.

\begin{figure}[t]
  \centering
  \begin{tikzpicture}[
    >=Stealth,
    layer/.style={draw, thick, rounded corners=3pt, minimum width=7cm,
      minimum height=1.1cm, font=\small, align=center},
    sublbl/.style={font=\tiny, text=black!60},
    arr/.style={->, thick, gray!70},
  ]
    % Layer 4 (bottom): Data Layer
    \node[layer, fill=green!12] (data) at (0,0)
      {\textbf{Data Layer}};
    \node[sublbl, below=-0.15cm of data]
      {JSON \quad Parquet/CSV \quad Schedule Resolution};

    % Layer 3: Model Layer
    \node[layer, fill=blue!10, above=0.5cm of data] (model)
      {\textbf{Model Layer}};
    \node[sublbl, below=-0.15cm of model]
      {Planning $\to$ System $\to$ Components};

    % Layer 2: LP Assembly Layer
    \node[layer, fill=orange!12, above=0.5cm of model] (lp)
      {\textbf{LP Assembly Layer}};
    \node[sublbl, below=-0.15cm of lp]
      {PlanningLP $\to$ SimulationLP $\to$ Component LPs};

    % Layer 1 (top): Solver Layer
    \node[layer, fill=red!10, above=0.5cm of lp] (solver)
      {\textbf{Solver Layer}};
    \node[sublbl, below=-0.15cm of solver]
      {CLP/CBC \quad CPLEX \quad HiGHS \quad (pluggable)};

    % Arrows between layers
    \draw[arr] (data.north) -- (model.south);
    \draw[arr] (model.north) -- (lp.south);
    \draw[arr] (lp.north) -- (solver.south);

    % Side labels
    \node[font=\tiny, rotate=90, anchor=south, gray]
      at (-3.9,1.2) {Data flow};
    \draw[->, thin, gray] (-3.9,0) -- (-3.9,3.6);
  \end{tikzpicture}
  \caption{\gtopt{} layered architecture: data ingestion, model
    construction, LP assembly, and pluggable solver backends.}
  \label{fig:arch_layers}
\end{figure}

The software is organized into four layers:

\begin{enumerate}
  \item \textbf{Data Layer}: JSON parsing (via DAW JSON Link, a
    zero-allocation header-only library), Parquet/CSV I/O (via Apache
    Arrow C++ SDK), and schedule resolution (\texttt{FieldSched<T>}
    variant type supporting scalar, vector, or filename references).
  \item \textbf{Model Layer}: The \texttt{Planning} $\rightarrow$
    \texttt{System} $\rightarrow$ component hierarchy that maps JSON
    input to C++ data structures.
  \item \textbf{LP Assembly Layer}: \texttt{PlanningLP} orchestrates
    \texttt{SimulationLP} instances (one per scenario/stage), which in
    turn assemble component LP objects (\texttt{GeneratorLP},
    \texttt{DemandLP}, \texttt{LineLP}, \texttt{BatteryLP},
    \texttt{ReservoirLP}, etc.) into a shared
    \texttt{FlatLinearProblem}.
  \item \textbf{Solver Layer}: The \texttt{SolverBackend} interface
    abstracts LP/MIP solution.  Implementations are loaded as dynamic
    plugins from configurable search paths.
\end{enumerate}

\subsection{Data Model and JSON Interface}

The input is specified as one or more JSON files containing three sections:

\begin{lstlisting}[language={},caption={Top-level JSON structure}]
{
  "options": { ... },
  "simulation": {
    "block_array": [...],
    "stage_array": [...],
    "scenario_array": [...]
  },
  "system": {
    "bus_array": [...],
    "generator_array": [...],
    "demand_array": [...],
    "line_array": [...],
    "battery_array": [...],
    ...
  }
}
\end{lstlisting}

Time-series data (demand profiles, generator capacity factors, reservoir
inflows) can be specified inline as JSON arrays, or as references to
external Parquet or CSV files in the \texttt{input\_directory}.  This
hybrid approach keeps the JSON configuration compact while supporting
arbitrarily large time-series via columnar Parquet storage.

\subsection{Python Toolchain}

\gtopt{} includes a comprehensive Python toolchain (16 registered CLI
tools):

\begin{itemize}
  \item \textbf{igtopt}: Excel workbook ($\rightarrow$ JSON + Parquet)
    converter for spreadsheet-based case preparation.
  \item \textbf{plp2gtopt}: PLP/PROMAX ($\rightarrow$ gtopt) converter
    for legacy hydrothermal scheduling data.
  \item \textbf{pp2gtopt / gtopt2pp}: pandapower
    ($\leftrightarrow$ gtopt) bidirectional converters for power flow
    model exchange.
  \item \textbf{gtopt\_compare}: Automated cross-validation against
    pandapower DC OPF on IEEE test networks.
  \item \textbf{gtopt\_diagram}: Network topology diagram generator
    using Mermaid syntax.
  \item \textbf{run\_gtopt / sddp\_monitor}: Execution management and
    SDDP convergence monitoring.
  \item \textbf{Validation tools}: \texttt{gtopt\_check\_json},
    \texttt{gtopt\_check\_lp}, \texttt{gtopt\_check\_output},
    \texttt{gtopt\_check\_solvers} for input validation and diagnostics.
\end{itemize}

\subsection{Web and GUI Services}

Two companion services provide graphical interfaces:

\begin{itemize}
  \item \textbf{webservice}: A Next.js REST API and browser UI for
    submitting optimization jobs, monitoring progress, and browsing
    results.
  \item \textbf{guiservice}: A Python/Flask application for creating
    and editing cases through a web-based form interface.
\end{itemize}
